% AnacostiaIQ Sensor REST API — reference
% Build:  pdflatex rest-api.tex   (twice, for the table of contents)
\documentclass[11pt,a4paper]{article}

\usepackage[margin=2.4cm]{geometry}
\usepackage[T1]{fontenc}
\usepackage[utf8]{inputenc}
\usepackage{lmodern}
\usepackage{xcolor}
\usepackage{listings}
\usepackage{booktabs}
\usepackage{longtable}
\usepackage{enumitem}
\usepackage[hidelinks]{hyperref}

\definecolor{codebg}{HTML}{F6F6F4}
\definecolor{codekw}{HTML}{2A78D6}
\definecolor{codestr}{HTML}{1BAF7A}
\definecolor{codecom}{HTML}{52514E}
\definecolor{warnbg}{HTML}{FDF3E7}

\lstdefinestyle{shell}{
  basicstyle=\ttfamily\footnotesize,
  backgroundcolor=\color{codebg},
  frame=single, framerule=0pt, framesep=6pt,
  breaklines=true, breakatwhitespace=false,
  postbreak=\mbox{\textcolor{codecom}{$\hookrightarrow$}\space},
  columns=fullflexible, keepspaces=true,
  showstringspaces=false,
  commentstyle=\color{codecom}\itshape,
}
\lstdefinestyle{js}{
  style=shell,
  language=Java,
  keywordstyle=\color{codekw}\bfseries,
  stringstyle=\color{codestr},
  morekeywords={const,let,async,await,function,return,for,of,if,else,new,try,catch,throw},
}
\lstset{style=shell}

\newcommand{\http}[2]{\texttt{\textbf{#1}}~\texttt{#2}}

\title{\textbf{AnacostiaIQ Sensor REST API}\\[2mm]
       \large Reference and command syntax}
\author{}
\date{Generated 6 October 2026}

\begin{document}
\maketitle
\tableofcontents
\bigskip

%===============================================================
\section{Overview}

A Flask service stores sensor readings in Amazon DynamoDB and serves them
over HTTP. Field stations \texttt{POST} readings; the dashboard and any
other client \texttt{GET} them.

\begin{center}
\begin{tabular}{@{}ll@{}}
\toprule
Base URL        & \texttt{http://54.213.147.59:5000} \\
Host            & \texttt{ec2-54-213-147-59.us-west-2.compute.amazonaws.com} \\
Service         & \texttt{sensor-api.service} (systemd), \texttt{\textasciitilde/sensor-api/app.py} \\
Storage         & DynamoDB table \texttt{sensor\_readings}, region \texttt{us-west-2} \\
Content type    & \texttt{application/json} \\
CORS            & enabled for all origins \\
\bottomrule
\end{tabular}
\end{center}

\subsection{Data model}

Each reading is one item:

\begin{center}
\begin{tabular}{@{}llp{7.6cm}@{}}
\toprule
Field & Type & Meaning \\
\midrule
\texttt{project}   & string & Namespace. Groups one deployment's stations. \\
\texttt{sensor\_id}& string & Sensor name, unique \emph{within} a project. \\
\texttt{timestamp} & string & ISO~8601, \texttt{YYYY-MM-DDTHH:MM:SS}. Sort key. \\
\texttt{value}     & string & Always stored and returned as a \emph{string}. \\
\texttt{unit}      & string & Free text, e.g. \texttt{m}, \texttt{\%}, \texttt{mm}. \\
\bottomrule
\end{tabular}
\end{center}

\subsection{How projects work}

\label{sec:projects}
The DynamoDB partition key is \texttt{sensor\_id} and the sort key is
\texttt{timestamp}, so a given pair identifies exactly one item. Without
namespacing, two projects posting the same sensor name at the same instant
would overwrite each other silently.

The API therefore stores the partition key as
\texttt{<project>\#<sensor\_id>} and splits it apart again on read. Clients
only ever see \texttt{project} and \texttt{sensor\_id} as separate fields;
the \texttt{\#} encoding never appears on the wire.

\begin{itemize}[leftmargin=*]
  \item \texttt{project} must not contain \texttt{\#}. A request that
        breaks this is rejected with \texttt{400}.
  \item Omitting \texttt{project} reads and writes the \emph{unnamespaced}
        space, which is what the API did before projects existed.
  \item \texttt{value} is stored by \texttt{str()}, so \texttt{42.5} comes
        back as \texttt{"42.5"}. Cast on the client.
  \item A \texttt{POST} to an existing \texttt{(project, sensor\_id,
        timestamp)} \textbf{overwrites} it. That is the only way to correct
        a reading --- there is no \texttt{PUT} or \texttt{PATCH}.
\end{itemize}

%===============================================================
\section{Endpoints}

\begin{center}
\begin{tabular}{@{}lll@{}}
\toprule
Method & Path & Purpose \\
\midrule
\texttt{GET}    & \texttt{/health}            & liveness check \\
\texttt{POST}   & \texttt{/sensor}            & store one reading \\
\texttt{GET}    & \texttt{/sensor/<id>}       & read a sensor's readings \\
\texttt{DELETE} & \texttt{/sensor/<id>}       & delete readings (token required) \\
\texttt{GET}    & \texttt{/sensors}           & list sensor ids \\
\texttt{GET}    & \texttt{/projects}          & list project names \\
\bottomrule
\end{tabular}
\end{center}

\noindent There is no authentication on reads or writes. Only
\texttt{DELETE} is gated, by the \texttt{X-Wipe-Token} header.

%---------------------------------------------------------------
\subsection{\http{GET}{/health}}

Returns \texttt{\{"status":"ok"\}} if the service is up.

\begin{lstlisting}
curl http://54.213.147.59:5000/health
\end{lstlisting}

%---------------------------------------------------------------
\subsection{\http{POST}{/sensor}}

Stores one reading. Body fields:

\begin{center}
\begin{tabular}{@{}lll@{}}
\toprule
Field & Required & Default \\
\midrule
\texttt{sensor\_id} & yes & --- \\
\texttt{value}      & yes & --- \\
\texttt{project}    & no  & none (unnamespaced) \\
\texttt{unit}       & no  & \texttt{""} \\
\texttt{timestamp}  & no  & server UTC time \\
\bottomrule
\end{tabular}
\end{center}

\begin{lstlisting}
curl -X POST http://54.213.147.59:5000/sensor \
     -H 'Content-Type: application/json' \
     -d '{"project":"John-McCormack",
          "sensor_id":"lab_moisture_a0",
          "value":42.5,
          "unit":"%",
          "timestamp":"2026-10-06T19:00:00"}'
\end{lstlisting}

\noindent Letting the server supply the timestamp:

\begin{lstlisting}
curl -X POST http://54.213.147.59:5000/sensor \
     -H 'Content-Type: application/json' \
     -d '{"project":"John-McCormack","sensor_id":"lab_moisture_a0","value":42.5,"unit":"%"}'
\end{lstlisting}

\noindent Responses: \texttt{200} \texttt{\{"status":"success"\}};
\texttt{400} if \texttt{project} contains \texttt{\#}.

%---------------------------------------------------------------
\subsection{\http{GET}{/sensor/<sensor\_id>}}

Returns an array of readings, oldest first.

\begin{center}
\begin{tabular}{@{}lp{10.4cm}@{}}
\toprule
Query parameter & Meaning \\
\midrule
\texttt{project} & restrict to one project; omit for unnamespaced data \\
\texttt{start}   & inclusive lower bound on \texttt{timestamp} \\
\texttt{end}     & inclusive upper bound on \texttt{timestamp} \\
\bottomrule
\end{tabular}
\end{center}

\noindent \texttt{start} and \texttt{end} are compared as \emph{strings}
against the ISO timestamp, so a prefix such as \texttt{2026-10} works, but
the format must match the stored one. Either may be given alone.

\begin{lstlisting}
# everything for one sensor in one project
curl "http://54.213.147.59:5000/sensor/lab_moisture_a0?project=John-McCormack"

# a date range
curl "http://54.213.147.59:5000/sensor/lab_moisture_a0?project=John-McCormack&start=2026-10-01&end=2026-10-07"

# everything since a moment
curl "http://54.213.147.59:5000/sensor/lab_moisture_a0?project=John-McCormack&start=2026-10-06T00:00:00"

# pretty-printed
curl -s "http://54.213.147.59:5000/sensor/precip_amount?project=John-McCormack" | python3 -m json.tool
\end{lstlisting}

\noindent Example response:

\begin{lstlisting}
[{"project":"John-McCormack","sensor_id":"lab_moisture_a0",
  "timestamp":"2026-10-06T19:00:00","unit":"%","value":"42.5"}]
\end{lstlisting}

%---------------------------------------------------------------
\subsection{\http{GET}{/sensors}}

Lists sensor ids. With \texttt{?project=} only that project's sensors;
without it, every sensor in every project.

\begin{lstlisting}
curl "http://54.213.147.59:5000/sensors?project=John-McCormack"
curl "http://54.213.147.59:5000/sensors"
\end{lstlisting}

\noindent This performs a full table scan, so it grows more expensive as
the table does.

%---------------------------------------------------------------
\subsection{\http{GET}{/projects}}

Lists distinct project names. A \texttt{null} entry means readings that
carry no project.

\begin{lstlisting}
curl http://54.213.147.59:5000/projects
\end{lstlisting}

%---------------------------------------------------------------
\subsection{\http{DELETE}{/sensor/<sensor\_id>}}

\begin{center}\fcolorbox{black}{warnbg}{\parbox{0.93\textwidth}{\textbf{Destructive.}
Requires the \texttt{X-Wipe-Token} header. Omitting both \texttt{start} and
\texttt{end} deletes \emph{every} reading for that sensor in that project.
Without a token configured server-side, deletion is refused outright.}}\end{center}

\medskip
Accepts the same \texttt{project}, \texttt{start} and \texttt{end}
parameters as the \texttt{GET}, and returns the number removed.

\begin{lstlisting}
# one sensor, one project, a bounded range
curl -X DELETE "http://54.213.147.59:5000/sensor/lab_moisture_a0?project=John-McCormack&start=2026-10-01&end=2026-10-02" \
     -H "X-Wipe-Token: $WIPE_TOKEN"

# every reading for that sensor in that project
curl -X DELETE "http://54.213.147.59:5000/sensor/lab_moisture_a0?project=John-McCormack" \
     -H "X-Wipe-Token: $WIPE_TOKEN"
\end{lstlisting}

\noindent Responses: \texttt{200}
\texttt{\{"status":"success","sensor\_id":...,"deleted":N\}};
\texttt{401} \texttt{\{"error":"unauthorized"\}}.

The token lives in the service environment on the server:

\begin{lstlisting}
sudo systemctl show -p Environment --value sensor-api
\end{lstlisting}

%===============================================================
\section{JavaScript}

All examples use \texttt{fetch}, and work in a browser or in Node~18+.

\subsection{Posting a reading}
\begin{lstlisting}[style=js]
const API = "http://54.213.147.59:5000";
const PROJECT = "John-McCormack";

async function postReading(sensorId, value, unit) {
  const res = await fetch(`${API}/sensor`, {
    method: "POST",
    headers: { "Content-Type": "application/json" },
    body: JSON.stringify({
      project: PROJECT,
      sensor_id: sensorId,
      value: value,
      unit: unit,
      timestamp: new Date().toISOString().slice(0, 19)  // no trailing Z
    })
  });
  if (!res.ok) throw new Error(`POST failed: ${res.status}`);
  return res.json();
}
\end{lstlisting}

\subsection{Reading a range}
\begin{lstlisting}[style=js]
async function getReadings(sensorId, start, end) {
  const q = new URLSearchParams({ project: PROJECT });
  if (start) q.set("start", start);
  if (end)   q.set("end", end);

  const res = await fetch(`${API}/sensor/${encodeURIComponent(sensorId)}?${q}`);
  if (!res.ok) throw new Error(`GET failed: ${res.status}`);

  const rows = await res.json();
  // value arrives as a string -- cast before plotting
  return rows.map(r => ({ t: new Date(r.timestamp), v: parseFloat(r.value) }));
}

const series = await getReadings("lab_moisture_a0",
                                 "2026-10-01", "2026-10-07");
\end{lstlisting}

\subsection{Listing sensors and projects}
\begin{lstlisting}[style=js]
const sensors  = await (await fetch(`${API}/sensors?project=${PROJECT}`)).json();
const projects = await (await fetch(`${API}/projects`)).json();
// projects may contain null, meaning readings with no project
\end{lstlisting}

\subsection{Deleting (token required)}
\begin{lstlisting}[style=js]
async function deleteReadings(sensorId, token, start, end) {
  const q = new URLSearchParams({ project: PROJECT });
  if (start) q.set("start", start);
  if (end)   q.set("end", end);

  const res = await fetch(`${API}/sensor/${encodeURIComponent(sensorId)}?${q}`, {
    method: "DELETE",
    headers: { "X-Wipe-Token": token }
  });
  if (res.status === 401) throw new Error("unauthorized -- bad or missing token");
  return res.json();   // { status, sensor_id, project, deleted }
}
\end{lstlisting}

%===============================================================
\section{Operations}

\subsection{Server access}
\begin{lstlisting}
ssh -i ~/Dropbox/AWS_/ArashLinux.pem ubuntu@54.213.147.59
\end{lstlisting}

\subsection{Service control}
\begin{lstlisting}
systemctl status sensor-api
sudo systemctl restart sensor-api
journalctl -u sensor-api -f
\end{lstlisting}

\subsection{Backing up the table}
\begin{lstlisting}
~/sensor-api/venv/bin/python3 - <<'PY'
import boto3, gzip, json, datetime
t = boto3.resource("dynamodb", region_name="us-west-2").Table("sensor_readings")
path = "backup-%s.json.gz" % datetime.datetime.now(datetime.UTC).strftime("%Y%m%d-%H%M%S")
kw, n = {}, 0
with gzip.open(path, "wt") as fh:
    while True:
        r = t.scan(**kw)
        for it in r["Items"]:
            fh.write(json.dumps({k: str(v) for k, v in it.items()}) + "\n"); n += 1
        if "LastEvaluatedKey" not in r: break
        kw["ExclusiveStartKey"] = r["LastEvaluatedKey"]
print(path, n, "items")
PY
\end{lstlisting}

\subsection{Acceptance tests}
\begin{lstlisting}
python3 scripts/test_project_api.py
python3 scripts/test_project_api.py --token "$WIPE_TOKEN"   # adds DELETE tests
\end{lstlisting}

%===============================================================
\section{Client configuration}

The \texttt{project} value must agree on both sides, or the dashboard
charts nothing.

\begin{center}
\begin{tabular}{@{}lll@{}}
\toprule
Side & File & Key \\
\midrule
Station (Pi) & \texttt{config.json}          & \texttt{"project"} (top level) \\
Dashboard    & \texttt{FrontEnd/config.json} & \texttt{"project"} \\
\bottomrule
\end{tabular}
\end{center}

\begin{lstlisting}
{
  "project": "John-McCormack",
  "station": { "id": "field_station_01", "name": "John McCormack Field Station" },
  "app": { "apiUrl": "http://54.213.147.59:5000/sensor" }
}
\end{lstlisting}

\noindent \texttt{project} groups stations; \texttt{station.id} still
distinguishes them within it.

%===============================================================
\section{Known limitations}

\begin{enumerate}[leftmargin=*]
  \item \textbf{No authentication on reads or writes.} Port 5000 is open to
        the world and only \texttt{DELETE} checks a token. Anyone who can
        reach the host can post arbitrary readings or read everything.
  \item \textbf{No update endpoint.} Correcting a value means re-posting
        the same key, which overwrites.
  \item \textbf{\texttt{value} is a string.} Clients must cast.
  \item \textbf{\texttt{/sensors} and \texttt{/projects} scan the table.}
        Cost grows with table size; cache the result rather than polling.
  \item \textbf{Timestamps are compared lexicographically.} Mixing formats
        (with and without a trailing \texttt{Z}) breaks range queries.
        Post \texttt{YYYY-MM-DDTHH:MM:SS}.
  \item \textbf{Forecast and observation share the table.} Weather sensors
        such as \texttt{precip\_amount} carry future timestamps.
\end{enumerate}

\end{document}
